\chapter{$\zeta$ 共形映射、Hardy 轨迹与 $\pi_0$ 逼近}\label{ch:visualization}

上一章给出 $\pi_0$ 的数值逼近。本章\textbf{不引入新定理}，而用三组动画的关键帧，
把前文章节已建立的同一套对象画成三种几何：
\begin{enumerate}
  \item \textbf{共形映射}（§\ref{sec:conformal_global}，图~\ref{fig:zeta_conformal_still}）：
        $w=\zeta(s)$ 下的全局索引——网格、临界线、零点、极点同框。
  \item \textbf{Hardy 轨迹}（§\ref{sec:spiral}，图~\ref{fig:hardy_z_spiral_still}）：
        $s=\frac12+it$ 的值域路径；过原点 $\Leftrightarrow$ 临界线零点。
  \item \textbf{$\pi_0$ 逼近}（§\ref{sec:explicit_vis}，图~\ref{fig:pi0_approx_start}--\ref{fig:pi0_approx_end}）：
        同一组 $\rho_k$ 合成 $\pi_0(x)$，与 $\pi(x)$ 阶梯对照。
\end{enumerate}
三者顺序是：定义域总图 $\to$ 临界线切片 $\to$ 算术轴合成；
分别交叉引用第~\ref{ch:riemann_zeta_intro}--\ref{ch:zero_distribution}~章、
第~\ref{ch:riemann_numerical_approximation}~章的 $Z$ 与相位，
以及第~\ref{ch:riemann_explicit_formula}--\ref{ch:riemann_numerical_approximation}~章的显式公式。
正文只放截图；完整视频、Manim / Mathematica 源程序与运行环境说明见\textbf{书后附注}
（\url{https://github.com/lichengtong123/lichengtong-RH}）。

% ============================================================
\section{共形映射：\texorpdfstring{$\zeta(s)$}{Zeta(s)} 的全局几何索引}
\label{sec:conformal_global}
% ============================================================

在面向公众的黎曼猜想介绍中，$w=\zeta(s)$ 的网格变形（常称「保角变换」「\(\zeta\) 映射动画」）出现频率很高，
几乎成为该主题的一种视觉标志：画面曲线运动迅速复杂、具有强烈的视觉与概念冲击。
例如著名数学科普频道 3Blue1Brown（3B1B）关于 $\zeta$ 与解析延拓的专题中，
就有以类圆曲线扩张、发散为主的映射演示；但通常没有同时表现：
临界带与临界线在映射下的扭曲、相交，以及非平凡零点沿射线向值域原点 $w=0$ 的收缩。
而把临界带、临界线、零点与 $s=1$ 极点等几何上重要的元素留在旁白或略去，
作为数学普及无可厚非。

本节借用上述画面语言，把书中已建立的对象标进同一变形过程：
网格对应延拓后的 $\zeta$，红点对应零点（同伦下向 $w=0$ 收缩，并与 RH 的几何说法对照），
中央发亮对应 $s=1$ 极点与素数主项的同一奇点。
完整视频见书后附注；下文截图与文字对照前文章节。

\subsection{设置与截图}

动画对定义域网格作同伦
\begin{equation}
  f_\alpha(s)=(1-\alpha)s+\alpha\,\zeta(s),\qquad \alpha\in[0,1],
  \label{eq:homotopy}
\end{equation}
$\alpha=0$ 为恒等，$\alpha=1$ 为 $w=\zeta(s)$。
图~\ref{fig:zeta_conformal_still} 取变换进行中的一帧（画面约 $56\%$ 进度），
已足以读出标记；完整变形见书后附注中的 \texttt{ZetaConformalMap.mp4}。

\begin{figure}[htbp]
\centering
\includegraphics[width=0.92\textwidth]{全书图形/保角变换.pdf}
\caption{$\zeta$ 共形映射动画关键帧。
黄线为网格在 $f_\alpha$ 下的像；中央亮区对应 $s=1$ 邻域被推离有限平面；
竖直绿带为临界带示意，红点为临界线上的非平凡零点（及其共轭）。
完整动画见书后附注。}
\label{fig:zeta_conformal_still}
\end{figure}

\subsection{画面与前文章节的对应}

\paragraph{（1）网格：延拓后的 $\zeta$，不是半平面级数}
终态必须用\textbf{解析延拓后}的 $\zeta$（第~\ref{ch:analytic_continuation_general}~、
\ref{ch:riemann_zeta_continuation}~章）。
仅在 $\operatorname{Re}(s)>1$ 的 Dirichlet 级数无法合法铺满临界带内的网格；
画面能进入 $\operatorname{Re}(s)\le 1$ 并在零点、极点附近剧烈形变，
正是“延拓已经完成”的几何回声。
密网格窗有限（虚部约 $|t|\lesssim$ 数个单位），更高零点靠临界线外延标出——
这是示意截断，不是临界带只到 $|t|\sim 1$。

\paragraph{（2）两个公式：定义区与函数方程}
与网格对照，宜分清下列两式的适用范围。

\textbf{定义区}（第~\ref{ch:riemann_zeta_intro}~章；仅 $\operatorname{Re}s>1$）
\[
\zeta(s)=\sum_{n=1}^{\infty}n^{-s}
=\prod_{p}\bigl(1-p^{-s}\bigr)^{-1}
\quad(\operatorname{Re}s>1).
\]
它说明 $\zeta$ 在绝对收敛半平面上的来源（级数 / Euler 积），
但\textbf{不能}单独支撑本图临界带内的网格。

\textbf{函数方程}（第~\ref{ch:riemann_zeta_continuation}~章；与~\eqref{eq:asymmetric_fe_alt} 同型）
\[
\zeta(s)=2^s\pi^{s-1}\sin\!\Bigl(\frac{\pi s}{2}\Bigr)\Gamma(1-s)\,\zeta(1-s).
\]
它给出延拓后的全局关系（$\zeta(s)$ 与 $\zeta(1-s)$ 由 $\chi(s)$ 相联；
对称形式 $\Lambda(s)=\Lambda(1-s)$、$\xi(s)=\xi(1-s)$ 见第~\ref{ch:xi_function}~章），
\textbf{并不}依赖左右半平面网格的镜面演示。
简言之：级数管定义区，函数方程管延拓后的全局约束；网格所实现的是延拓后的 $\zeta$。

\paragraph{（3）红点：零点定义与 RH 的几何说法}
红点取 $\rho=\frac12\pm it_k$（数值表与第~\ref{ch:zero_distribution}~章 $N(T)$ 所用同一组对象）。
因 $\zeta(\rho)=0$，同伦下
$f_\alpha(\rho)=(1-\alpha)\rho$，标记沿射线收向原点 $w=0$——
这是零点定义在同伦下的直接推论。
书中已证：非平凡零点必在临界带 $0<\operatorname{Re}s<1$ 内
（第~\ref{ch:zero_distribution}~章）；
\textbf{黎曼猜想（RH）}的几何说法是：这些“坍缩中心”还须落在临界线
$\operatorname{Re}s=\frac12$ 的像上。
图中零点取在临界线上，是已有数值结果与 RH 工作假设下的惯用画法，\textbf{不能代替 RH 的证明}。
共轭成对红点对应 $\zeta(\overline s)=\overline{\zeta(s)}$（系数实）。

\paragraph{（4）中央发亮：极点 $s=1$ 与素数主项}
$\zeta$ 在 $s=1$ 的单极点（Laurent 主项 $1/(s-1)$，第~\ref{ch:riemann_zeta_intro}~章）
使小圆周像的半径 $\sim 1/r\to\infty$，网格被推离有限平面。
同一奇点在第~\ref{ch:riemann_explicit_formula}~章显式公式中贡献
$\psi$ 的主项 $x$ 或 $J$ 的 $\li(x)$——
值域侧“推向无穷”与素数计数的光滑主趋势是同一极点的两面。

\paragraph{（5）保角性与零点、极点}
第~\ref{ch:complex_analysis_intro}~章：全纯映射在 $f'\neq 0$ 处局部保角（夹角保持）。
对固定 $\alpha$，$f_\alpha'=(1-\alpha)+\alpha\zeta'$，
故仅在 $f_\alpha'\neq 0$ 的开集上局部正交，\textbf{不能}由画面断言全程处处保角。
对终态 $w=\zeta(s)$：
\begin{itemize}
  \item 极点 $s=1$：$\zeta$ 无定义，无保角性；画面中央发亮对应网格被推向无穷；
  \item 非平凡零点 $\rho$：$\zeta(\rho)=0$，保角性取决于 $\zeta'(\rho)$——
        若为单零点则 $\zeta'(\rho)\neq 0$，局部仍保角（网格交角 $90^\circ$）；
        若 $\zeta'(\rho)=0$（多重零点），该点不保角。
        数值上已知大量零点为单零点，\textbf{一般单零性未证}；
  \item 其余 $\zeta'\neq 0$ 处：局部保角，是画面中网格正交的数学依据。
\end{itemize}
共轭对称 $\zeta(\bar s)=\overline{\zeta(s)}$
（第~\ref{ch:riemann_zeta_intro}~章第~\ref{sec:schwarz}~节，用第~\ref{ch:complex_analysis_intro}~章 Schwarz 反射）
使上下半平面的像关于 $w$ 平面实轴对称，与画面上下成对标记一致。

\paragraph{阅读边界}
同伦~\eqref{eq:homotopy} 不是第三种延拓；证明仍在第~\ref{ch:riemann_zeta_continuation}~章。
本图红线（临界线）$\to$ 下节螺旋；本图红点 $\rho_k$ $\to$ 第三节 $\pi_0$ 的频率。

% ============================================================
\section{Hardy 轨迹：\texorpdfstring{$Z$}{Z} 与过原点}
\label{sec:spiral}
% ============================================================

\subsection{设置与截图}

取共形图中的临界线 $s=\frac12+it$，画值域路径
\[
\vec r(t)
=\Bigl(
  \operatorname{Re}\zeta\bigl(\tfrac12+it\bigr),\;
  \operatorname{Im}\zeta\bigl(\tfrac12+it\bigr),\;
  t
\Bigr).
\]
图~\ref{fig:hardy_z_spiral_still} 为配套动画关键帧。
临界线不是任选竖线：函数方程 $\Lambda(s)=\Lambda(1-s)$（或 $\xi(s)=\xi(1-s)$）
以 $\operatorname{Re}s=\frac12$ 为对称轴
（第~\ref{ch:riemann_zeta_continuation}~、\ref{ch:xi_function}~章），
故本节在函数方程所取的对称轴 $\operatorname{Re}s=\frac12$ 上观察 $\zeta$。

\begin{figure}[htbp]
\centering
\includegraphics[width=0.92\textwidth]{全书图形/Hardy Z函数.pdf}
\caption{Hardy 轨迹动画关键帧。
白线：三维路径 $\vec r(t)$；绿线：在 $(\operatorname{Re}\zeta,\operatorname{Im}\zeta)$ 平面上的投影；
橙 / 蓝：侧面投影 $\operatorname{Re}\zeta(\frac12+it)$、$\operatorname{Im}\zeta(\frac12+it)$。
白点 $N_k$ 标临界线零点；底栏为 Hardy 极坐标~\eqref{eq:hardy_polar}。
完整动画见书后附注。}
\label{fig:hardy_z_spiral_still}
\end{figure}

\subsection{画面与前文章节的对应}

\paragraph{（1）三维轨迹与第~\ref{ch:zero_distribution}~章的曲面}
白螺旋是 $\vec r(t)$；绿线是它在值域复平面上的投影。
侧面橙、蓝曲线
\[
t\mapsto\operatorname{Re}\zeta\bigl(\tfrac12+it\bigr),\qquad
t\mapsto\operatorname{Im}\zeta\bigl(\tfrac12+it\bigr)
\]
恰是第~\ref{ch:zero_distribution}~章实部 / 虚部高度曲面
在 $\sigma=\frac12$ 上的截线。
零迹线 $A=\{\operatorname{Re}\zeta=0\}$、$B=\{\operatorname{Im}\zeta=0\}$ 的交点为非平凡零点；
沿临界线仅当两投影在\textbf{同一} $t$ 同时过零，三维轨迹才过值域原点。

\paragraph{（2）过原点与 Hardy $Z$（第~\ref{ch:riemann_numerical_approximation}~章）}
底栏公式即该章的极坐标分解~\eqref{eq:hardy_polar}：
\[
\zeta\Bigl(\tfrac12+it\Bigr)
=Z(t)\,e^{-i\vartheta(t)},
\qquad
Z(t)=e^{i\vartheta(t)}\zeta\Bigl(\tfrac12+it\Bigr)\in\mathbb{R}.
\]
实值性来自函数方程与 $\vartheta$ 的定义（该章正文短证）；
$Z(\gamma)=0$ $\Leftrightarrow$ $\zeta(\frac12+i\gamma)=0$ $\Leftrightarrow$ 轨迹过原点。
第~\ref{ch:riemann_numerical_approximation}~章用 Gram / 割线求零点，
就是在找 $Z$ 的实根；本节是同一分解的几何显示。
模长 $|Z(t)|$ 决定螺圈半径，相位 $-\vartheta(t)$（及 $Z$ 的符号）决定旋转。

\paragraph{（3）螺距变密与 $N(T)$（第~\ref{ch:gamma_function}~、\ref{ch:zero_distribution}~章）}
Riemann--Siegel 相位的 Stirling 展开给出
$\vartheta'(t)\sim\frac12\log(t/2\pi)$，
故角速度随 $t$ 对数增长，螺旋在 $t$ 轴上\textbf{越绕越密}。
这与 Riemann--von Mangoldt 公式
\[
N(T)=\frac{T}{2\pi}\log\frac{T}{2\pi e}+O(\log T)
\]
的主项同源：平均零点间距 $\sim 2\pi/\log(t/2\pi)\to 0$。
画面上的“变距”= 书中零点密度增大的几何说法，不是独立现象。

\paragraph{（4）圈半径涨落}
$|Z(t)|$ 在相邻零点间的局部极大决定各圈半径；
形式 Dirichlet 级数在临界线上不绝对收敛，但相位 $t\log n$ 的叠加
提示几乎周期型干涉（与第~\ref{ch:riemann_numerical_approximation}~章
Riemann--Siegel 主项同一思想层次）。
“杂乱变幅”来自多频率，非单一谐波。

\paragraph{（5）与 RH、无零点区的关系}
若 RH 成立，全部非平凡零点都表现为本螺旋的过原点事件。
若存在 $\beta\neq\frac12$ 的零点，它\textbf{不会}出现在本曲线上
（而在竖直线 $\operatorname{Re}s=\beta$ 的值域轨迹上）。
$\sigma>1$ 上 Euler 乘积保证 $\zeta\neq 0$；$\sigma=1,\,t\neq 0$ 上无零点与素数定理等价
（第~\ref{ch:prime_distribution}~、\ref{ch:zero_distribution}~章）——
那些竖直线\textbf{不过}原点，与本图的“反复过原点”形成对照。

\paragraph{阅读边界}
曲线用延拓后的 $\zeta(\frac12+it)$，不是临界线上的绝对收敛级数。
画面只显示有限段 $t$；$N(T)$ 描述 $T\to\infty$ 的平均行为。
过原点标记的是\textbf{临界线上}的零点；RH 的全称断言不能由图代替。

\paragraph{与前后节}
过原点高度 $\{\gamma_k\}$ $=$ 共形图红点虚部 $=$ 下节 $\pi_0$ 振荡的频率表。

% ============================================================
\section{$\pi_0$ 逼近：从零点到素数阶梯}
\label{sec:explicit_vis}
% ============================================================

\subsection{设置与截图}

前两节把非平凡零点固定为：共形映射中坍向 $w=0$ 的点，
以及螺旋上 $t=\gamma_k$ 的过原点事件。
本节把\textbf{同一组} $\rho_k=\frac12+i\gamma_k$ 代入
第~\ref{ch:riemann_numerical_approximation}~章的 Möbius--$\Ei$ 截断~\eqref{eq:pi_explicit_full}，
在 $x$ 轴上画 $\pi_0$，并与 $\pi(x)$ 对照。
动画变量主要是非平凡零点个数 $K$；
图~\ref{fig:pi0_approx_start}、\ref{fig:pi0_approx_end}
为开始帧（$K=31$）与结束帧（$K=500$），平凡零点均取 $M=10$。

\begin{figure}[htbp]
\centering
\includegraphics[width=0.92\textwidth]{全书图形/黎曼逼近1.pdf}
\caption{$\pi_0$ 逼近动画开始帧：非平凡零点 $K=31$，平凡零点 $M=10$。
橙线 $\pi(x)$，青线为 $\pi_0$ 逼近；曲线较圆滑，棱角未充分形成。
完整动画见书后附注。}
\label{fig:pi0_approx_start}
\end{figure}

\begin{figure}[htbp]
\centering
\includegraphics[width=0.92\textwidth]{全书图形/黎曼逼近2.pdf}
\caption{同一动画结束帧：$K=500$，$M=10$。
$\pi_0$ 更贴阶梯平台，跳跃由振荡重构；有限零点逼近的波纹向素数间断点收缩
（与第~\ref{ch:riemann_numerical_approximation}~章图~\ref{fig:pi_comparison_merged} 同旨；
勿与 Fourier 级数的经典 Gibbs 常数混称）。
完整动画见书后附注。}
\label{fig:pi0_approx_end}
\end{figure}

\subsection{画面与前文章节的对应}

\paragraph{（1）橙 / 青两条曲线}
橙色 $\pi(x)=\#\{p\le x\}$ 是第~\ref{ch:elementary_number_theoretic_functions}~章的算术对象。
青色是~\eqref{eq:pi_explicit_full} 的有限零点逼近 $\pi_0^{(K,M)}$，
不是另造阶梯，而是第~\ref{ch:riemann_explicit_formula}~章显式公式对 $\pi$ 的数值拟合。

\paragraph{（2）公式链路：$\psi\to J\to\pi$}
第~\ref{ch:riemann_explicit_formula}~章主线先有
\[
\psi_0(x)=x-\sum_{\rho}\frac{x^{\rho}}{\rho}-\cdots,
\]
再经分部得 $J_0$，最后 Möbius 反演
\[
\pi(x)=\sum_{k\ge 1}\frac{\mu(k)}{k}\,J(x^{1/k})
\]
回到 $\pi$（第~\ref{ch:elementary_number_theoretic_functions}~章反演与 §\ref{sec:psi-to-pi}）。
图中出现的 $J$ 式、$\mu$ 定义与 $\pi_0$ 的三重 $\Ei$ 和，
即该链路在计算层的展开：
\begin{itemize}
  \item \textbf{主项} $\Ei(\log x)/n$ 型：来自极点 $s=1$——
        与共形图中央“推离有限平面”同一奇点；
  \item \textbf{非平凡项} $\Ei(\rho_k\log x/n)$：
        $\rho_k$ 即前两节红点 / 螺旋过原点高度；
  \item \textbf{平凡项} $\Ei((-2m/n)\log x)$：负偶数零点，随 $m$ 快衰（$M=10$ 已足）。
\end{itemize}
实现用 $\Ei$ 统一复宗量分支；$\Li$ 仅用于 $\pi\sim\Li$ 的误差比较
（见第~\ref{sec:Li}~节与第~\ref{ch:riemann_numerical_approximation}~章）。

\paragraph{（3）$K$ 增大：与第~\ref{ch:fourier_vs_riemann}~、\ref{ch:riemann_numerical_approximation}~章}
更多 $\gamma_k$ 叠加 $\Rightarrow$ 平台更平、跳跃更陡。
第~\ref{ch:fourier_vs_riemann}~章比较了 Fourier 整数频率与对数余弦 / 零点频率；
此处是后者在 $\pi_0$ 上的动态版。
有限 $K$ 下跳跃附近可有过冲与波纹，波纹可向间断点收缩——
与第~\ref{ch:riemann_numerical_approximation}~章图~\ref{fig:pi_comparison_merged}
（$K=20,60,200,500$ 静态对照）同一思想；逼近核与 Dirichlet 核不同，
过冲比例\textbf{不必}等于经典 Gibbs 约 $9\%$（第~\ref{ch:fourier_vs_riemann}~章）。
评价宜看\textbf{整体波形}；素数点处有限项逼近收敛于中值 $\pi(p)-\frac12$，
青线一般不等于 $\pi(p)$。

\paragraph{（4）振幅、$\beta_k$ 与 RH}
零点项振幅含因子 $x^{\beta_k}$（对 $\psi$ 约为 $x^{\beta}/|\rho|$）。
此处取 $\beta_k=\frac12$（临界线零点表），即在 RH 工作假设下运行；
若某 $\beta>\frac12$，该项会在更大尺度上主导（第~\ref{ch:prime_distribution}~章
von Koch 型误差与第~\ref{ch:riemann_explicit_formula}~章振幅分析）。
波形随 $K$ 改善只说明有限零点逼近有效，\textbf{不能}由此证明 RH 或素数定理。

\paragraph{阅读边界}
\eqref{eq:pi_explicit_full} 在有限 $K,M,N$ 时为近似；
完整等式需对称求和并处理条件收敛（第~\ref{ch:riemann_explicit_formula}~章）。
窗 $x\le 200$ 仅便于观看；大 $x$ 需更大 $K$ 与更稳的 $\Ei$ 求值。

% ============================================================
\section{三条示意的逻辑串联}
% ============================================================

\begin{center}
共形映射（$\rho$ 在定义域）
$\;\longrightarrow\;$
Hardy 轨迹（$t=\gamma$ 过原点）
$\;\longrightarrow\;$
$\pi_0$ 逼近（$\gamma$ 为阶梯频率）
\end{center}

\begin{itemize}
  \item \textbf{零点}：映射中的红点 $\to$ 轨迹过原点的高度 $\to$ 素数阶梯振荡频率
        （同一 $\{\rho_k\}$ 的三种角色）。
  \item \textbf{极点 $s=1$}：共形映射中推离有限平面 $\to$ $\pi_0$ 的光滑主项。
  \item \textbf{函数方程/$\xi$}：保证临界线为对称轴、零点可作 Hadamard / 显式公式骨架
        （第~\ref{ch:xi_function}~、\ref{ch:riemann_explicit_formula}~章已证）。
\end{itemize}
叙事与 1859 年论文一致：先 $\zeta$ 的全局与零点，再进入素数计数；
本章三条示意按这一逻辑顺序衔接。
全文知识谱系底图见第~\ref{ch:spectrum_map}~章图~\ref{fig:function-relations}。

\section*{本章小结与后续展望}
\addcontentsline{toc}{section}{本章小结与后续展望}

本章以三组几何示意对照前文章节：
$\zeta$ 共形映射（全局几何索引）、Hardy 轨迹（临界线上 $Z\,e^{-i\vartheta}$ 与 $N(T)$）、
$\pi_0$ 逼近（算术轴上的阶梯拟合）。
每节交叉引用前文章节的定义、函数方程、零点理论与数值公式；
证明与误差分析仍以第~\ref{ch:riemann_zeta_continuation}--\ref{ch:riemann_numerical_approximation}~章为准。
完整动画与源码见\textbf{书后附注}及 GitHub 仓库 \url{https://github.com/lichengtong123/lichengtong-RH}。

\paragraph{后续展望}
下一章解读 Riemann 1859 年原稿（第~\ref{ch:riemann_original_interpretation}~章），
在数值与几何工具齐备之后回到历史源头；再后为当代前沿简介
（第~\ref{ch:contemporary_research}~章）。
